\subsection{Orthogonale Projektion}
\label{subsec:vector_algebra_orthogonal_projection}

Ein Vektor $\vec b$ lässt sich relativ zu einer Richtung $\vec a\ne\vec 0$
in eine parallele und eine senkrechte Komponente zerlegen:
\[
  \boxed{\vec b=\vec b_{\parallel}+\vec b_{\perp}}.
\]

\begin{center}
  \CachedTikzFigure{content/resources/chapters/vector_algebra/tikz/orthogonal_projection.tex}
\end{center}

Mit dem Einheitsvektor $\vec e_a=\vec a/|\vec a|$ ergibt sich die
vorzeichenbehaftete Komponente von $\vec b$ in Richtung $\vec a$:
\begin{ClozeArea}[34]{5}
  \[
    \begin{aligned}
      \vec b\cdot\vec e_a
      &=(\vec b_{\parallel}+\vec b_{\perp})\cdot\vec e_a\\[0.5em]
      &=\vec b_{\parallel}\cdot\vec e_a
        +\underbrace{\vec b_{\perp}\cdot\vec e_a}_{=0}.
    \end{aligned}
  \]
\end{ClozeArea}

\begin{KeyPointBox}{Parallele und senkrechte Komponente}{vector_algebra_projection_components}
  \[
    \boxed{
      \vec b_{\parallel}
      =(\vec b\cdot\vec e_a)\,\vec e_a
      =\frac{\vec a\cdot\vec b}{\vec a\cdot\vec a}\,\vec a
    },
  \]
  \[
    \boxed{
      \vec b_{\perp}
      =\vec b-\vec b_{\parallel}
    }.
  \]
  Dabei gilt $|\vec b_{\parallel}|=|\vec b\cdot\vec e_a|$ und
  $\vec a\cdot\vec b_{\perp}=0$.
\end{KeyPointBox}

\begin{ExampleBox}{Einen Vektor orthogonal zerlegen}{vector_algebra_projection_example}
  Gegeben seien
  \[
    \vec a=\begin{pmatrix}2\\1\end{pmatrix},
    \qquad
    \vec b=\begin{pmatrix}1\\3\end{pmatrix}.
  \]
  Dann gilt
  \[
    \vec b_{\parallel}
    =\frac{\vec a\cdot\vec b}{\vec a\cdot\vec a}\,\vec a
    =\frac5{5}\begin{pmatrix}2\\1\end{pmatrix}
    =\Cloze{2.5cm}{\begin{pmatrix}2\\1\end{pmatrix}},
  \]
  \[
    \vec b_{\perp}
    =\vec b-\vec b_{\parallel}
    =\Cloze{2.5cm}{\begin{pmatrix}-1\\2\end{pmatrix}},
    \qquad
    \vec a\cdot\vec b_{\perp}=0.
  \]
\end{ExampleBox}
